\documentclass[10pt]{article}
\usepackage[margin=20mm]{geometry}
\usepackage{amsmath,amssymb,mathtools,fontspec,unicode-math,xurl,hyperref}
\IfFontExistsTF{DejaVu Serif}{\setmainfont{DejaVu Serif}}{\setmainfont{Latin Modern Roman}}
\setmathfont{latinmodern-math.otf}
\setlength{\parindent}{0pt}\setlength{\parskip}{6pt}\sloppy
\title{Seifert surfaces in the four-ball: diagrammatic and algebraic checks, a genus correction, and review limits}
\author{AI-assisted technical review}
\date{}
\begin{document}
\maketitle
Abstract: The printed Seifert matrices give the claimed integral intersection-form obstruction, and the local Frobenius-algebra identities used to distinguish stabilizations hold over the specified field of two elements. One Minor error is confirmed: Remark 4.4 gives the wrong genus after adding a tube. No Major error has been established. In this continuation, all 33 main pages and all five supplementary pages were actually inspected, with enlarged checks of the Khovanov movies, framed links, torus-surgery moves, and Stein diagram. These checks support the pictured local arguments. The author’s computations were not executed, and a complete independent reconstruction of the Seiberg–Witten cap and basic-class argument was not obtained. The conclusion remains limited to the coverage stated below and is not an independent acceptance decision.\par
\section*{Exact text and equivalence relations}
The bound paper is Kyle Hayden, Seungwon Kim, Maggie Miller, JungHwan Park, and Isaac Sundberg, Seifert surfaces in the 4-ball, arXiv:2205.15283v3, uploaded on 2 May 2023 at 02:41:43 UTC [paper]. The supplied main PDF has 33 pages; the associated Additional documentation has five pages [supp]. Both complete texts, including references and the displayed software transcripts, were read. The arXiv record describes the v3 change as a correction to the floating of Figure 27. The journal association is DOI 10.4171/JEMS/1703; the publisher lists acceptance on 18 April 2025 and online publication on 29 September 2025 [journal]. Identity with that later published text was not established.\par
Three distinctions are essential throughout. A boundary-fixed isotopy is stronger than an ambient isotopy allowing motion on the boundary. A topological obstruction is stronger than a smooth obstruction. Minimal Seifert genus is asserted for some examples, whereas the Whitehead-double surfaces used for the simplest cobordism-map distinction need not be minimal. The report does not transfer a claim from one of these settings to another.\par
Theorem 1.3 and the intersection forms provide topological nonisotopy for the basic pair. For the Whitehead doubles, Alexander polynomial one and the cited topological classification supply topological isotopy relative to the boundary. The Khovanov maps then give a smooth distinction. Upgrading a boundary-fixed distinction to an unrestricted ambient distinction requires the separate symmetry argument or the asymmetric-knot construction; it does not follow solely from different maps for a chosen diagram.\par
\section*{Independent algebraic checks}
\subsection*{The orientable branched covers}
On page 21, the printed matrices are \[V_0=\begin{pmatrix}-6&-2\\-3&-2\end{pmatrix},
 \qquad
 V_1=\begin{pmatrix}-6&0\\1&-1\end{pmatrix}.\] Exact symbolic calculation gives \[\det(V_i-tV_i^T)=6t^2-11t+6,\qquad
 Q_0=V_0+V_0^T=\begin{pmatrix}-12&-5\\-5&-4\end{pmatrix},
 \qquad
 Q_1=\begin{pmatrix}-12&1\\1&-2\end{pmatrix}.\] Both forms have determinant 23 and are negative definite. The Alexander polynomial has span two, consistent with minimal genus one for the displayed genus-one surfaces. Figure 18 was subsequently inspected at enlarged resolution. Its framed links have diagonal framings \((-12,-4)\) and \((-12,-2)\); the signed mixed-crossing sums are \(-10\) and \(2\), giving off-diagonal linking numbers \(-5\) and \(1\). The shown basis curves intersect once. The doubled-band construction, together with this skew intersection, recovers the printed matrices, up to compatible basis orientations.\par
Equation reference: eq:square. The vector \((0,1)\) has square \(-2\) under \(Q_1\). Write \(q(m,n)=6m^2+5mn+2n^2\). Then \[
 8q(m,n)=(4n+5m)^2+23m^2.\] If \(q=1\), this identity forces \(m=0\), and then \(2n^2=1\) is impossible. The equation \(q=-1\) is impossible by nonnegativity. Thus \(Q_0\) represents neither \(2\) nor \(-2\). This includes the orientation-reversing possibility for a proposed homeomorphism of the branched covers. It verifies the integer obstruction in Theorem 4.1. The separate actual comparison of the coloured basis curves, doubled bands, and framed links supplies the previously unchecked diagrammatic input. The general branched-cover algorithm is used as the cited standard result; it is not reproved here.\par
\subsection*{The nonorientable variant and the rank-one case}
The forms printed in Remark 4.5 are \[R_0=\begin{pmatrix}-12&-5\\-5&-3\end{pmatrix},
 \qquad
 R_1=\begin{pmatrix}-12&1\\1&-1\end{pmatrix}.\] They have determinant 11. The second represents \(-1\); the first does not, because \[12(12m^2+10mn+3n^2)=(6n+10m)^2+44m^2.\] For a value \(1\), this forces \(m=0\), followed by \(3n^2=1\); a negative value is again excluded.\par
For Corollary 4.6, a one-minimum ribbon Möbius band gives the stated single-two-handle branched-cover description. Figure 21 was actually inspected: the disk and nonorientable band lift to the displayed single trefoil handle, with framing \(9\) in that example. In the general one-minimum case the form is \([\pm\det J]\). When \(\det J>1\), no integer vector has square \(1\) or \(-1\), excluding a projective-plane summand by the indicated branched-cover argument. The pretzel determinant computation also checks: \[\left|(2n)(-2n-1)+(-2n-1)(2n+1)+(2n)(2n+1)\right|
 =(2n+1)^2\quad(n\geq1).\]\par
\subsection*{What the Frobenius algebra does and does not verify}
Over \(\mathbb F_2\), take \(A=\mathbb F_2[x]/(x^2)\) with \[\Delta(1)=1\otimes x+x\otimes1,\quad
 \Delta(x)=x\otimes x,\quad
 \epsilon(1)=0,\quad\epsilon(x)=1.\] Then \(\mu(x\otimes x)=0\) and \(\mu\Delta=0\) on both basis vectors. These exact identities support the algebraic mechanisms in Propositions 3.1 and 3.5: merging two \(x\)-labeled components kills the generator, and an added handle kills the cobordism map. The calculation was performed independently as matrices over \(\mathbb F_2\).\par
The geometric input was checked separately in the continuation. All five rows of Figure 11 and the three local models in Figure 12 were actually inspected at enlarged resolution. A further 400-dpi crop makes the gray \(0\)-resolution bands of the initial state legible: each joins distinct resolved components. Thus every outgoing differential on the all-\(x\) state vanishes. Following the pictured saddles and Reidemeister moves keeps the all-\(x\) term with coefficient one and ends in disk boundaries; each final counit is \(\epsilon(x)=1\). The complementary surface admits a saddle merging distinct \(x\)-components and hence gives zero. This checks the specific state and movie using the standard cited Reidemeister chain maps; no Khovanov program was run.\par
For the Whitehead double, the three Figure 12 models have compatible paired-sheet resolutions, and the clasp removes an \(x\)-labelled circle with coefficient one. Figure 13 supplies the complementary zero-producing saddle. The doubled quasipositive models in Figure 15 and the trefoil extension in Figure 17 were also traced. In particular, the nonzero movie value is not inferred solely from the abstract identity \(\mu\Delta=0\). That identity separately shows why a stabilized surface has zero map over \(\mathbb F_2\).\par
\section*{Confirmed finding}
Problem 1 (Minor: genus after stabilization).  Page 22, Remark 4.4, first sentence, describes adding a tube to each of the genus-one, one-boundary-component surfaces \(\Sigma_0,\Sigma_1\) and calls the resulting surfaces genus one. The genus must be two.\par
For either original surface, \(\chi(\Sigma_i)=2-2(1)-1=-1\). An interior tube removes two disks and inserts an annulus. Consequently the new Euler characteristic is \(-3\), and the boundary still has one component: \[-3=2-2g-1,\qquad g=2.\] This calculation uses exactly the operation stated in the remark. It is not a conjectural objection about the drawing. The same wording is present in the official version-specific HTML.\par
The correction is to replace genus-1 by genus-2 in this sentence. The statement about stabilization distance one is consistent with genus two and need not be changed on this account. Figure 19 has now been inspected. Its intermediate band movie allows the endpoints to slide on the surgered link, after which the other three bands are restored. This is consistent with the stated relative-boundary isotopy in the four-ball. The correction remains a local genus misstatement, not evidence against the main nonisotopy theorems.\par
\section*{Proof dependencies and coverage limits}
\subsection*{Satellites, symmetry, and minimal genus}
The genus comparison in Proposition 2.5 has the required strict inequality: removing two copies of a single-boundary-component subsurface of genus \(f\geq1\) and reconnecting by an annulus changes the genus by \(-2f+1<0\). The single boundary component matters for the connectivity and Euler-characteristic calculation.\par
Proposition 2.6 supplies the special meridional-torus-twist isotopy for the satellites in question. It is not a statement that all torus twists of arbitrary Seifert surfaces become boundary-fixed isotopies. Proposition 2.9 analyzes the relevant knot symmetries and is used to extend boundary motion over the surface pair. Figures 8–10 were actually inspected: the outer parallel sheets can be untwisted through the product region in the four-ball, and the pictured half-turn representatives preserve the two surface pieces. The exhaustion of all boundary symmetries still uses the cited JSJ classification. The extension across a collar in Proposition 2.10 is compatible with a time reparametrization constant near the inner end of the collar.\par
The strongly quasipositive construction in Theorem 1.6 distinguishes a non-destabilizable surface from a stabilized one. This is an intrinsic distinction and avoids identifying maps after an arbitrary boundary symmetry. The specified genus is \(2g(J)\); minimal genus is not asserted for that pair.\par
\subsection*{The ancillary computations}
The five-page supplement contains substantial supporting information: Dowker–Thistlethwaite input, a genus computation, a hyperbolic symmetry argument, and a finite-slope loop. It is therefore incorrect to describe the paper as having no computational support. All five supplementary pages, including Figures A–C and every displayed command, were actually inspected. The crossing circle in Figure C surrounds the connecting band; its \((1,-n)\) filling has the stated full-twist interpretation. None of the author’s commands was executed during this review.\par
Proposition A.1 uses the computed genus-three base case together with the band-sum genus inequality and its equality case. The argument for the family depends on that base case, not merely on the existence of a genus-three Seifert surface.\par
In Section B, the systole bound, the normalized-length threshold, and the finite list of exceptional slopes play different roles. Using the printed cusp area \(18.8600\ldots\), the numerical ratio \(46/\sqrt{18.86}\) is about \(10.59\), still larger than \(10.1\). The text’s rounded use of \(18.6\) therefore does not defeat the required numerical separation. This arithmetic is conditional on the actual cusp-area and slope computations.\par
The supplement explicitly uses `verified=True` for the short-slope and canonical-retriangulation computations, while displaying ordinary length-spectrum and cusp-area calls for the preliminary bounds. Current official documentation distinguishes floating-point output from verified output, and notes that `verify\_hyperbolicity` returns a success Boolean [snappy]. The supplement does not supply a complete versioned transcript with interval certificates for every intermediate bound. This review does not establish that any reported geometric value is false, nor does current documentation alone establish the behavior of the exact historical installation. The genus base case and the asymmetry computations remain unreproduced dependencies.\par
\subsection*{Infinite families and gauge theory}
For the disconnected family in Section 4.2, the relative-boundary hypothesis permits attaching the same band before comparing branched covers. The resulting first-homology distinction is between \(\mathbb Z\) and \(\mathbb Z_{2|n|}\) when \(n\neq0\). To compare two parameters, one first reduces to their difference by undoing a common torus twist. These are different assertions from free ambient isotopy of the original surfaces in the three-sphere.\par
The actual Figure 23 attaching curves also give this homology calculation directly: in the one-generator abelianization the red two-handle imposes \(2n\) times the generator, while the gray handle imposes zero. The framings are not substituted for these attaching exponents.\par
For the connected family, Figures 24–28 were actually compared. In Figure 26 the initial form, with a compatible orientation choice, is \(\left(\begin{smallmatrix}-4&4\\4&-4\end{smallmatrix}\right)\). Sliding to the sum of its two generators gives the displayed zero-framed component. The two cancelling handle pairs open the twist boxes into the marked product-torus neighbourhood. In Figure 27 the two parallel surgeries carry the same shear \(\gamma\mapsto\gamma+n\beta\); their composition sends \(\gamma\) to \(\gamma+2n\beta\). An independent exact matrix calculation confirms this coefficient for every integer \(n\), and the last isolated zero-framed circle cancels the added three-handle. The boundary markings were followed in the actual diagram, not replaced by this arithmetic check.\par
In the final front of Figure 28, the red and blue components each have three cusp pairs, the gray component has one, and the orange component has none in the displayed one-handle convention. The corresponding Thurston–Bennequin values are \(-3,-3,-1,0\); the printed framings \(-4,-4,-2,-1\) are one less. The added \(-1\) sphere meets the square-zero torus once, so smoothing their intersection gives a genus-one class of square \(-1+0+2=1\), as required by the adjunction step.\par
These checks resolve the previously unviewed local handle comparisons. The subsequent Stein embedding and Morgan–Mrowka–Szabó formula remain external foundations. The printed formula controls sums over translates by the torus class. A complete independent reconstruction of the asserted support of individual basic classes and all cap-induced homology markings was not obtained. Cancellation cannot be excluded by arithmetic alone, and signed parameters require the stated relative-boundary identifications. No actual counterexample to that part of the argument was found. These are explicit limits on the review, not confirmed Major findings.\par
\section*{Supporting checks}
The independent script verifies the two Alexander determinants, the two integral square identities, determinants and represented vectors, the \(\mathbb F_2\) multiplication/comultiplication identities, the tube Euler characteristic, the pretzel polynomial identity, and the stated rounded cusp-threshold inequality. A supplementary finite check covers 40,401 integer pairs; the universal nonrepresentation conclusions follow from the displayed square identities, not from that finite search.\par
The first execution stopped on the reviewer’s use of structural equality between expanded and factored symbolic expressions. The failed script, exit status, and output were retained. The corrected script checks the polynomial difference and exits successfully. This was a reviewer-script issue and is not a manuscript finding.\par
The earlier image-access failures, including the original `EMFILE` records and false inspection flags, are preserved. The continuation used the documented native image output to display the exact allocated PNG bytes in this scientific conversation. All main and supplementary pages were actually viewed, followed by the enlarged geometric checks described above. A separate reviewer-written script checks the handle-slide form, composition of the two surgery shears, Stein framing arithmetic, and the sphere–torus square; it exits successfully and executes no author code.\par
The previous report and associated records are retained. This revised TeX is genuinely recompiled with the approved Tectonic runtime and the paper’s own retained dependencies. Rendering alone does not count as viewing. The supported assessment is one Minor finding and no established Major, within the limits above. Scientific decision remains pending and human reviews remain empty; the separately recorded delivery state denotes a completed initial-review package, not independent acceptance.\par
9 K. Hayden, S. Kim, M. Miller, J. Park, and I. Sundberg, Seifert surfaces in the 4-ball, arXiv:2205.15283v3 (2023), https://arxiv.org/abs/2205.15283v3. The same authors, Additional documentation for Seifert surfaces in the 4-ball, five-page v3 ancillary PDF, https://arxiv.org/src/2205.15283v3/anc/additional\_documentation.pdf. EMS Press, article record for DOI 10.4171/JEMS/1703, https://ems.press/journals/jems/articles/14299167. SnapPy developers, Verified computations and Manifold: the main class, official documentation consulted 27 October 2026, https://snappy.computop.org/verify.html, https://snappy.computop.org/manifold.html.\par
\end{document}
